% =====================================================================
%  Trabalho de Graduação (TG) — FATEC Indaiatuba / ADS
%  Hierarquia de normas: Manual de Normas FATEC-ID (2026-2027) > ABNT.
%  Compilar com: make   (pdflatex + bibtex + pdflatex x2)
% =====================================================================
\documentclass[
  12pt,
  oneside,
  a4paper,
  sumario=tradicional,
  brazil
]{abntex2}

% ----------------------------- Pacotes -------------------------------
\usepackage[T1]{fontenc}
\usepackage[utf8]{inputenc}
\usepackage{mathptmx}            % Times New Roman (texto e matemática)
\usepackage[brazil]{babel}
\usepackage{graphicx}
\usepackage{array}
\usepackage{booktabs}
\usepackage{multirow}
\usepackage{longtable}
\usepackage{amsmath,amssymb}
\usepackage{xcolor}
\usepackage{listings}
\usepackage{microtype}
\usepackage[alf]{abntex2cite}    % citações ABNT autor-data (estilo alf)

\graphicspath{{graficos/}}

% --- Tipos de ilustração (Manual FATEC-ID, item 3) --------------------
% Quadro = dado qualitativo; Tabela = dado numérico; Gráfico = representação
% geométrica de dados; Figura = o resto. O abnTeX2 só traz figura e tabela.
% Quadro e gráfico escrevem no mesmo arquivo da lista de figuras (.lof):
% como têm menos de 5 itens cada, o Manual (4.1.1.10) manda fazer uma lista
% única, chamada LISTA DE ILUSTRAÇÕES. A lista de tabelas é separada.
% Modelos do Manual (3.1 e 3.2): tabela com laterais abertas e linhas entre as
% linhas; quadro fechado, em grade. Um pouco de altura nas linhas para as
% réguas não encostarem no texto.
\renewcommand{\arraystretch}{1.2}
\newfloat{quadro}{lof}{Quadro}
\newlistentry{quadro}{lof}{0}
\renewcommand{\cftquadroname}{Quadro\space}
\renewcommand*{\cftquadroaftersnum}{\hfill\textendash\hfill}
\newfloat{grafico}{lof}{Gráfico}
\newlistentry{grafico}{lof}{0}
\renewcommand{\cftgraficoname}{Gráfico\space}
\renewcommand*{\cftgraficoaftersnum}{\hfill\textendash\hfill}
% Posição das ilustrações (22/10/2026, pedido do autor): sem posição fixa. Todas
% usam [htbp] e ficam o mais próximo possível do trecho que as cita, na mesma
% página ou no topo da seguinte (Manual FATEC-ID, item 3). O [section] do
% placeins impede que uma ilustração atravesse o título da seção seguinte, e o
% cap3 tem \FloatBarrier antes de cada "O ano de". A regra antiga de figuras
% fixas ([H], pacote float, de 15/10) foi retirada; cópia em
% notas/rascunhos/antes_imagens_2026-10-22/.
\usepackage[section]{placeins}
\renewcommand{\topfraction}{0.9}
\renewcommand{\bottomfraction}{0.8}
\renewcommand{\textfraction}{0.1}
\renewcommand{\floatpagefraction}{0.75}
\setcounter{topnumber}{3}
\setcounter{totalnumber}{4}
% Página só de ilustrações: empilhadas a partir do topo, como numa página com
% texto, em vez de espalhadas com vãos entre elas.
\makeatletter
\setlength{\@fptop}{0pt}
\setlength{\@fpsep}{\floatsep}
\setlength{\@fpbot}{0pt plus 1fil}
\makeatother

% ----------------------- Ajustes de formatação -----------------------
% Manual FATEC: recuo de primeira linha 1,25 cm; espaçamento 1,5;
% margens 3/2/3/2 (padrão abnTeX2).
\setlength{\parindent}{1.25cm}
\setlength{\parskip}{0pt}
\setlength{\baselineskip}{1.5\baselineskip}

% Desabilitar remoção automática de parindent no primeiro parágrafo de seções
\usepackage{indentfirst}

% --- Notas de rodapé (ABNT NBR 6028 + Manual FATEC 1.3 e 1.4) ---
% Conforme ABNT: linha próxima da nota, afastada do texto
\renewcommand{\footnotesize}{\fontsize{10}{12}\selectfont}
\setlength{\footnotesep}{0.3cm}  % espaço entre texto e linha
\renewcommand{\footnoterule}{\hrule height 0.4pt width 3cm depth 0pt \vskip 2pt}  % linha com 3cm, próxima da nota

% --- Estilo de títulos (Manual FATEC 1.3 e 4.2.2) --------------------
% Capítulo: 14 pt, negrito, só inicial maiúscula, alinhado à esquerda.
% Subtítulos (seção/subseção): 12 pt, negrito, alinhados à esquerda.
\renewcommand{\chaptitlefont}{\normalfont\bfseries\fontsize{14}{18}\selectfont}
\renewcommand{\chapnumfont}{\chaptitlefont}
% O abnTeX2 redefine \chaptitlefont no \AtBeginDocument; reaplicamos DEPOIS
% para garantir que o tamanho dos títulos seja realmente respeitado.
% Manual FATEC 1.3: títulos = 14 pt. (Para aproximar do PDF original, use 12.)
\AtBeginDocument{%
  \renewcommand{\chaptitlefont}{\normalfont\bfseries\fontsize{14}{18}\selectfont}%
  \renewcommand{\chapnumfont}{\chaptitlefont}%
}

% Rótulo de divisão centralizado (INTRODUÇÃO, CAPÍTULO I, ...).
\newif\ifhasdiv
\newcommand{\divlabel}{}
\newcommand{\setdiv}[1]{\hasdivtrue\renewcommand{\divlabel}{#1}}
\newif\ifchapnumbered

% Estilo de capítulo:
%  - capítulo numerado  -> "N Título" à esquerda, só inicial maiúscula;
%  - capítulo sem número (pré-textual/REFERÊNCIAS) -> centralizado, CAIXA-ALTA;
%  - rótulo de divisão (INTRODUÇÃO/CAPÍTULO) centralizado acima do título.
\makechapterstyle{fatec}{%
  \setlength{\beforechapskip}{0pt}
  \setlength{\afterchapskip}{\onelineskip}
  \setlength{\midchapskip}{0pt}
  \renewcommand{\chapterheadstart}{%
    \chapnumberedfalse
    \ifhasdiv
      \thispagestyle{empty}%  abertura de capítulo: contada, sem número
      \phantomsection
      \addcontentsline{toc}{chapter}{\divlabel}%
      {\centering\chaptitlefont\divlabel\par}%
      \vskip\onelineskip
      \hasdivfalse
    \fi
  }%
  \renewcommand{\printchaptername}{}%
  \renewcommand{\chapternamenum}{}%
  \renewcommand{\printchapternum}{\chapnumberedtrue\chaptitlefont\thechapter\hspace{0.5em}}%
  \renewcommand{\afterchapternum}{}%
  \renewcommand{\printchaptertitle}[1]{%
    \ifchapnumbered
      \chaptitlefont ##1%
    \else
      {\centering\chaptitlefont\MakeTextUppercase{##1}\par}%
    \fi}%
  \renewcommand{\afterchaptertitle}{\par\nobreak\vskip\onelineskip}%
}
\chapterstyle{fatec}

% Paginação: apenas o número no canto superior direito, sem cabeçalho
% corrente nem filete (Manual FATEC 1.5).
\makepagestyle{numtopright}
\makeevenhead{numtopright}{\thepage}{}{}
\makeoddhead{numtopright}{}{}{\thepage}
\makeheadrule{numtopright}{0pt}{0pt}
\makeevenfoot{numtopright}{}{}{}
\makeoddfoot{numtopright}{}{}{}
\aliaspagestyle{chapter}{numtopright}

% Títulos pré-textuais (RESUMO, SUMÁRIO, LISTA DE..., REFERÊNCIAS):
% abnTeX2 usa, por padrão, sans-serif \Huge — o Manual FATEC pede Times 14.
\renewcommand{\ABNTEXchapterfont}{\normalfont\bfseries\fontsize{14}{18}\selectfont}
\renewcommand{\ABNTEXchapterfontsize}{\fontsize{14}{18}\selectfont}
\renewcommand{\resumoname}{RESUMO}
\renewcommand{\bibname}{REFERÊNCIAS}
\renewcommand{\apendicesname}{APÊNDICES}

% Seções e subseções: 12 pt negrito, à esquerda (sem caixa-alta).
\setsecheadstyle{\normalfont\bfseries\fontsize{12}{16}\selectfont\raggedright}
\setsubsecheadstyle{\normalfont\bfseries\fontsize{12}{16}\selectfont\raggedright}
\setsubsubsecheadstyle{\normalfont\bfseries\fontsize{12}{16}\selectfont\raggedright}

% --- Sumário: 2 níveis (capítulo e seção); 1.1.1 omitido; todas as
%     entradas à esquerda (sem identação progressiva) e em negrito;
%     pontilhado denso em todos os níveis (Manual 1.5 + legibilidade).
\makeatletter
\renewcommand{\l@part}[2]{}                     % omite "APÊNDICES" do sumário
\makeatother
\renewcommand{\cftdotsep}{1}                    % pontilhado populado
\renewcommand{\cftchapterdotsep}{\cftdotsep}    % pontilhado também nos capítulos
\renewcommand{\cftsectiondotsep}{\cftdotsep}
\renewcommand*{\cftchapterfont}{\normalfont\bfseries}        % títulos em negrito
\renewcommand*{\cftsectionfont}{\normalfont}                 % subseções sem negrito
\renewcommand*{\cftchapterpagefont}{\normalfont\bfseries}
\renewcommand*{\cftsectionpagefont}{\normalfont}
\cftsetindents{chapter}{0em}{1.4em}
\cftsetindents{section}{0em}{2.6em}
\cftsetindents{subsection}{0em}{3.4em}         % Manual 4.1.1.14: indicativos à esquerda
\renewcommand*{\cftsubsectionfont}{\normalfont}
\renewcommand*{\cftsubsectionpagefont}{\normalfont}
\renewcommand{\cftsubsectiondotsep}{\cftdotsep}
\setlength{\cftbeforesubsectionskip}{0pt}
\setlength{\cftbeforepartskip}{0pt}            % apêndice omitido do sumário
% Reduzir espaçamento no sumário conforme ABNT (espaçamento simples)
\setlength{\cftbeforechapterskip}{0pt}         % sem espaço antes de capítulos
\setlength{\cftbeforesectionskip}{0pt}         % sem espaço antes de seções

% Garante nomes em caixa-alta no sumário (babel reescreve esses rótulos).
\addto\captionsbrazil{%
  \renewcommand{\bibname}{REFERÊNCIAS}%
  \renewcommand{\apendicesname}{APÊNDICES}%
  \renewcommand{\anexosname}{ANEXOS}%
  \renewcommand{\listfigurename}{LISTA DE ILUSTRAÇÕES}%
  \renewcommand{\listtablename}{LISTA DE TABELAS}%
}
\renewcommand{\listfigurename}{LISTA DE ILUSTRAÇÕES}
\renewcommand{\listtablename}{LISTA DE TABELAS}

% --- Legendas de ilustrações/tabelas (Manual FATEC 1.3 e seção 3) -----
% Tamanho 10, "Figura N" em negrito, travessão, legenda em cima e fonte
% logo abaixo da ilustração.
\renewcommand{\ABNTEXfontereduzida}{\fontsize{10}{12}\selectfont}
\captionnamefont{\fontsize{10}{12}\selectfont\bfseries}
\captiontitlefont{\fontsize{10}{12}\selectfont}
\captiondelim{\space--\space}
\captionstyle{\centering}
\setlength{\abovecaptionskip}{2pt}
\setlength{\belowcaptionskip}{2pt}
% A linha de fonte usa o comando nativo \fonte{...} do abnTeX2,
% que imprime "Fonte: ..." em tamanho reduzido.

% --- Estilo do código MQL5 (Apêndice A) ------------------------------
\definecolor{codegray}{gray}{0.4}
\definecolor{codegreen}{rgb}{0,0.5,0}
\definecolor{codeblue}{rgb}{0,0,0.6}
\lstdefinestyle{mql5}{
  language=C++,
  basicstyle=\ttfamily\fontsize{8}{9}\selectfont,
  keywordstyle=\color{codeblue}\bfseries,
  commentstyle=\color{codegreen},
  stringstyle=\color{codegray},
  numbers=left,
  numberstyle=\tiny\color{codegray},
  numbersep=8pt,
  breaklines=true,
  showstringspaces=false,
  tabsize=2,
  extendedchars=true,
  literate=%
    {á}{{\'a}}1 {é}{{\'e}}1 {í}{{\'i}}1 {ó}{{\'o}}1 {ú}{{\'u}}1
    {â}{{\^a}}1 {ê}{{\^e}}1 {ô}{{\^o}}1 {ã}{{\~a}}1 {õ}{{\~o}}1
    {ç}{{\c c}}1 {à}{{\`a}}1 {Á}{{\'A}}1 {É}{{\'E}}1 {Ç}{{\c C}}1
    {ô}{{\^o}}1 {Ã}{{\~A}}1 {Õ}{{\~O}}1 {â}{{\^a}}1,
  morekeywords={input,datetime,ulong,double,int,bool,string,void,
    OnInit,OnTick,OnDeinit,CTrade,MqlTick,ENUM_MA_METHOD,
    ENUM_APPLIED_PRICE,ENUM_POSITION_TYPE,iMA,CopyBuffer,
    SymbolInfoTick,SymbolInfoInteger,SymbolInfoDouble,NormalizeDouble,
    PositionsTotal,PositionGetTicket,PositionSelectByTicket,
    PositionGetString,PositionGetInteger,SeriesInfoInteger}
}

% Metadados PDF e navegação (links/marcadores não alteram a formatação
% impressa: colorlinks=false + pdfborder={0 0 0} deixam o texto igual ao
% padrão ABNT, só adicionam a navegação clicável).
\usepackage{xurl}           % URLs longas nas referências quebram em qualquer ponto
\usepackage{hyperref}
\hypersetup{
  pdftitle={Desenvolvimento e comparação de modelos de negociação algorítmica: ARIMA-GARCH e médias móveis},
  pdfauthor={Bruno Rocha Pereira e Giovanni Silva Loose},
  colorlinks=false, pdfborder={0 0 0}
}
\usepackage{bookmark}          % marcadores (bookmarks) do PDF mais confiáveis
\usepackage[all]{hypcap}       % links de figura/tabela apontam para o topo, não a legenda

% =====================================================================
